% =======================================================================================
\section{The Whole Band: B}
\label{app:b}
% =======================================================================================
% STATUS (2026-10-08): B1 is running on TREX (jobs 13882121 -> 13882126, 4 chained 24 h
% links + 1 spare). Epoch 91 of 200 at 07:45 UTC on 8 Oct, flow loss 0.526, 1536 s per
% epoch, no NaN; the cooldown starts at epoch 161, earliest finish ~08:30 CEST on 10 Oct.
% Everything marked [TBD] is filled from its scorecard and eval:
%   sbatch slurm/lisa-scorecard.sbatch B-late --n_random 50     # 60 windows, 240 sources, ~1 h
%   sbatch slurm/lisa-eval.sbatch B-late                         # corner pages, chunked draws
% then: copy scorecard.npz to paper/data/scorecards/B-late.npz, rerun paper/scripts.

\subsection{Setup}

B covers the whole band, 0.1--12\,mHz, in windows of 4096 bins ($R=1024$, an interior of 3072
bins) with 2048 data tokens, 64 times as many as XS. Everything else is the XS recipe
(\cref{app:xs}): four sources per window, the same priors, the 512-wide network (75.6M
parameters, unchanged), the warped clock with $p=3$, auxiliary heads for the first 10\% of the
steps, and a linear cooldown over the last 20\%. A step costs 1.54\,s on one A100 (41 times
XS) with a peak of 18.6\,GiB, so B can afford a fifth of the steps of the XS runs:

\begin{table}[h]
\caption{Cost of a B run on one A100, with the last 20\% of the steps cooling down, in 24-hour
jobs. At a fixed budget the 512-wide network gets 1.64 times the steps of the 768.}
\label{tab:b-cost}
\centering
\small
\begin{tabular}{lll}
\toprule
steps & width 512 & width 768 \\
\midrule
100k & 44\,h, 2 jobs & 71\,h, 3 jobs \\
150k & 65\,h, 3 jobs & 107\,h, 5 jobs \\
200k & 87\,h, 4 jobs & 143\,h, 6 jobs \\
\bottomrule
\end{tabular}
\end{table}

We train the 512-wide network for 200k steps, the last 40k cooling down. On XS, width 768
bought detection at equal steps (\cref{app:xs-ladder}), but B is far short of the XS step
counts, and each doubling of the steps was worth about 20\% in width on XS; at this budget,
steps are the better use of compute. \Cref{fig:training}c shows the run so far: the flow loss
fell from 0.753 to 0.526 in 91k steps and is still falling, with no instability. As on XS, the
auxiliary losses drift once their weight reaches zero (at 20k steps).

\subsection{Results}

\begin{table}[h]
\caption{B by SNR band, on 240 sources in 60 windows (\cref{app:evaluation}). \tbd}
\label{tab:b-band}
\centering
\small
\begin{tabular}{lrcccrccc}
\toprule
SNR $\snr$ & $n$ & found & width [bins] & $\times$ ideal & offset [bins] & in68 & in95 & rank \\
\midrule
$<15$ & \tbd & \tbd & \tbd & \tbd & \tbd & \tbd & \tbd & \tbd \\
15--40 & \tbd & \tbd & \tbd & \tbd & \tbd & \tbd & \tbd & \tbd \\
40--100 & \tbd & \tbd & \tbd & \tbd & \tbd & \tbd & \tbd & \tbd \\
$\geq 100$ & \tbd & \tbd & \tbd & \tbd & \tbd & \tbd & \tbd & \tbd \\
\midrule
all & \tbd & \tbd & & & & \tbd & \tbd & \tbd \\
\bottomrule
\end{tabular}
\end{table}

\begin{table}[h]
\caption{B by frequency, in thirds of the log band. Above a few mHz the chirp is resolved and
the chirp mass becomes measurable; the Doppler sideband widens and the sky position sharpens.
Chirp-mass and sky widths are for found sources at SNR $\ge40$. \tbd}
\label{tab:b-frequency}
\centering
\small
\begin{tabular}{lccccc}
\toprule
$f_0$ & found ($\snr\ge40$) & loud $f_0$ width [bins] & $\sigma_{\ln\Mc}$ & sky width [deg] & in95 \\
\midrule
0.1--0.49\,mHz & \tbd & \tbd & \tbd & \tbd & \tbd \\
0.49--2.4\,mHz & \tbd & \tbd & \tbd & \tbd & \tbd \\
2.4--12\,mHz & \tbd & \tbd & \tbd & \tbd & \tbd \\
\bottomrule
\end{tabular}
\end{table}

\begin{figure}[h]
  \centering
  \begin{minipage}{0.48\textwidth}\centering
    \placeholderfig[0.18\textheight]{$f_0$ width against SNR on B, with the XS final model and
    the ideal width (as \cref{fig:width-snr}).}
  \end{minipage}\hfill
  \begin{minipage}{0.48\textwidth}\centering
    \placeholderfig[0.18\textheight]{Corner plot of a loud high-frequency source, where the
    chirp mass is measured (as \cref{fig:corner-loudest}).}
  \end{minipage}
  \caption{B results. \tbd}
  \label{fig:b-results}
\end{figure}

The analysis that will fill these placeholders:
\begin{itemize}
  \item \textbf{Does the recipe carry over?} Detection, widths and calibration by SNR band
        (\cref{tab:b-band}), against the XS final model on the same SNR bands. \tbd
  \item \textbf{What the band adds.} At high frequency the chirp is resolved and the chirp
        mass measured (\cref{tab:b-frequency}); the Fisher forecast is then informative for
        $\Mc$ and is the reference for its width. \tbd
  \item \textbf{Does the wider window cost precision?} The B and XS models, compared on
        windows below 1.26\,mHz, where both are trained. \tbd
  \item \textbf{Resolution of the frequency coordinate.} One unit of $x_f$ spans 1536 bins on
        B, so its bfloat16 cell reaches 6 bins over half of the window
        (\cref{app:geometry-remarks}). The $f_0$ width against the source's position in the
        window, as in \cref{fig:bias}, tests whether this limits B. \tbd
\end{itemize}

% Predictions written before the result (2026-10-08), for the authors; not for the paper:
%  - the loss drops at the cooldown (epoch 161) as on XS, by ~0.02;
%  - if the recipe carries over: loud f0 floor near XS's 0.04 bins (25-30x ideal), SNR 15-40
%    detection >= 80%, in95 >= 0.95 in every band;
%  - if the bf16 cell of x_f matters on B: the loud width grows with |x_f| and sources with
%    |x_f| > 0.5 (half of them) are wider than a bin, so loud-source detection falls well
%    below XS's 99%. On XS the effect is at most 1.1-1.2x at a 0.094-bin cell; B's cell is
%    64x larger in bins. The fix would be fp32 Fourier features of the source coordinates
%    before the bf16 cast, which needs a new run.
